\documentclass[pdflatex,sn-mathphys-num]{sn-jnl}% Math and Physical Sciences Numbered Reference Style

\usepackage{graphicx}%
\usepackage{multirow}%
\usepackage{amsmath,amssymb,amsfonts}%
\usepackage{amsthm}%
\usepackage{mathrsfs}%
\usepackage[title]{appendix}%
\usepackage{xcolor}%
\usepackage{textcomp}%
\usepackage{manyfoot}%
\usepackage{booktabs}%
\usepackage{algorithm}%
\usepackage{algorithmicx}%
\usepackage{algpseudocode}%
\usepackage{listings}%
\usepackage{tikz}
\usetikzlibrary{shapes.geometric, arrows, positioning}

\theoremstyle{thmstyleone}%
\newtheorem{theorem}{Theorem}%  
\newtheorem{proposition}[theorem]{Proposition}% 

\theoremstyle{thmstyletwo}%
\newtheorem{example}{Example}%
\newtheorem{remark}{Remark}%

\theoremstyle{thmstylethree}%
\newtheorem{definition}{Definition}%

\raggedbottom

\begin{document}

\title[AdaptInfer: Distributed Microservices Architecture for LLM Strategy Orchestration]{AdaptInfer: A Scalable Distributed Microservices Architecture and Dynamic Strategy Orchestration Engine for Cost-Aware LLM Inference}

\author*[1]{\fnm{Adithya} \snm{R.}}\email{adithya.research@university.edu}
\author[1]{\fnm{Aniketh Reddy} \snm{J.}}\email{aniketh.reddy@university.edu}
\author[1]{\fnm{AmanReddy} \snm{J.}}\email{aman.reddy@university.edu}
\author[1]{\fnm{Naishadha} \snm{Badithala}}\email{naishadha.b@university.edu}

\affil*[1]{\orgdiv{Department of Computer Science and Engineering}, \orgname{Amrita School of Computing, Amrita Vishwa Vidyapeetham}, \orgaddress{\city{Chennai}, \postcode{601103}, \state{Tamil Nadu}, \country{India}}}

\abstract{\textbf{Abstract---}Production deployment of Large Language Models (LLMs) faces an acute systems engineering bottleneck: the severe trade-off between inference accuracy, token cost, and response latency. Statically assigning exhaustive multi-step reasoning strategies (e.g., Tree-of-Thoughts) across all incoming workloads incurs prohibitive median latencies ($45.2$\,s) and excessive token overhead ($1,959$ tokens/query), while static lightweight zero-shot prompting on small models suffers a catastrophic $46.1\%$ failure rate on complex multi-step reasoning. Existing LLM routers predominantly rely on rigid offline heuristics or binary model selection that fail to adapt dynamically to diverse query complexities and real-time hardware budget constraints. In this paper, we present \textbf{AdaptInfer}, an asynchronous distributed microservices architecture and dynamic strategy orchestration engine engineered for real-time joint selection of model parameter size ($3\text{B}$ vs. $7\text{B}$) and reasoning strategy ($\text{Direct}$, $\text{Chain-of-Thought}$, $\text{Graph-of-Thoughts}$, $\text{Tree-of-Thoughts}$) under session token budget constraints. AdaptInfer decouples synchronous local model generation (Ollama inference workers via Ray task scheduling) from asynchronous API state management (FastAPI microservice). The core orchestration engine dynamically evaluates a $390$-dimensional context vector comprising semantic sentence embeddings, structural complexity metrics, domain indicators, and remaining session budget ratios. To eliminate hallucination risks during online execution, AdaptInfer incorporates a \textit{Confidence-Gated Cascade} grounded in the classic Chain of Responsibility software design pattern, evaluating semantic answer consensus across dual sampling temperatures ($T=0.2$ vs. $T=0.8$) to escalate uncertain generations automatically. Evaluated across $3,712$ benchmark queries spanning mathematical, logical, and commonsense reasoning (GSM8K, LogiQA, CommonsenseQA), AdaptInfer achieves statistical accuracy parity with the theoretical Tree-of-Thoughts upper bound ($81.4\%$ vs. $82.0\%$) while slashing average response latency by \textbf{67.2\%} ($14.8$\,s vs. $45.2$\,s) and saving \textbf{27.5\% in token costs} ($1,421$ vs. $1,959$ tokens/query). Finally, we deliver an open-source, full-stack Single-Page Application (SPA) dashboard providing real-time telemetry, parallel strategy racing, and live performance visualization.}

\keywords{Large Language Models, Distributed Systems Architecture, Software Design Patterns, Strategy Pattern, Chain of Responsibility, Adaptive Orchestration, Real-Time Inference, AdaptInfer}

\maketitle

\section{Introduction}\label{sec1}

The rapid proliferation of Large Language Models (LLMs) across enterprise software systems has catalyzed a fundamental paradigm shift in automated cognitive processing. Advanced prompting mechanisms---such as Chain-of-Thought (CoT) \cite{wei2022chain}, Graph-of-Thoughts (GoT) \cite{besta2024graph}, and Tree-of-Thoughts (ToT) \cite{yao2023tree}---enable smaller models to execute complex logical, mathematical, and commonsense deduction by constructing explicit reasoning graphs. However, deploying these sophisticated reasoning frameworks in production environments exposes a fundamental systems engineering challenge: the severe tension along the \textit{Pareto Frontier} between inference latency, computational token expenditure, and output correctness.

\subsection{The Engineering Problem: Static Allocation Failure}
In production software architectures, statically assigning a single model or reasoning strategy across all incoming user traffic represents a major structural failure point:
\begin{enumerate}
    \item \textbf{Over-Provisioning Failure:} Directing every query to a high-capacity model (e.g., Qwen-7B) executing Tree-of-Thoughts (7B\_ToT) achieves maximum task accuracy (82.0\%). However, it imposes an extreme latency penalty (averaging 45.2 seconds per query) and massive token consumption (1,959 tokens per query), rapidly exceeding hardware hosting capacity and violating real-time Service Level Agreements (SLAs).
    \item \textbf{Under-Provisioning Failure:} Conversely, routing all traffic to a fast, low-parameter model executing zero-shot direct prompting (3B\_Direct) yields rapid response times (4.4 seconds), but results in an unacceptable 46.1\% error rate on multi-step reasoning problems.
\end{enumerate}

\subsection{Limitations of Existing Routing Solutions}
To resolve this trade-off, recent literature has explored LLM routing mechanisms. Frameworks such as RouteLLM \cite{ong2025routellm} and FrugalGPT \cite{chen2023frugalgpt} implement binary classifiers to route between strong and weak models based on human preference data. However, these systems suffer from three critical architectural limitations:
\begin{itemize}
    \item \textbf{Offline Training Rigidity:} They depend on offline supervised training loops, leaving them highly vulnerable to domain distribution shifts when exposed to live production traffic.
    \item \textbf{Stateless Budget Ignorance:} They treat model selection as an isolated decision, failing to adapt dynamically to depleting hardware session budgets or floating SLA thresholds.
    \item \textbf{Synchronous Thread Locking:} Existing implementations frequently couple routing evaluation with synchronous model inference, locking main application event loops during long-running generation.
\end{itemize}

\subsection{Our Solution: AdaptInfer Systems Architecture}
In this paper, we present \textbf{AdaptInfer}, a decoupled, asynchronous, distributed microservice architecture designed for real-time joint selection of model parameter size ($3\text{B}$ vs. $7\text{B}$) and reasoning strategy ($\text{Direct}$, $\text{CoT}$, $\text{GoT}$, $\text{ToT}$). AdaptInfer formulates query routing as a stateless decision problem governed by the LinUCB algorithm \cite{li2010contextual}. By embedding a real-time session budget scalar directly into a 384-dimensional Sentence-BERT feature space, the system mathematically encourages frugality as available compute depletes. 

To ensure high reliability during the decision engine's early online learning phases, we design a \textit{Confidence-Gated Cascade} implementing the Chain of Responsibility software design pattern. This gate evaluates logit distribution stability across sampling temperatures ($T=0.2$ vs. $T=0.8$), escalating ambiguous generations to higher-capacity fallback actions before output delivery.

\subsection{Key Technical Contributions}
The primary systems and software engineering contributions of this work are as follows:
\begin{enumerate}
    \item \textbf{Decoupled Distributed Architecture:} We engineer a distributed inference pipeline using the Ray computing framework and local Ollama instances, isolating synchronous LLM generation from asynchronous API state updates.
    \item \textbf{Object-Oriented Design Pattern Blueprint:} We implement a production-grade software design pattern foundation incorporating the \textit{Strategy Pattern} for polymorphic strategy execution, the \textit{Chain of Responsibility Pattern} for confidence-gated escalation, and the \textit{Adapter Pattern} for local LLM client isolation.
    \item \textbf{State-Augmented Dynamic Budget Tracker:} We concatenate a real-time remaining budget scalar $b_t$ directly to the dense semantic state vector $\phi(x_t)$, enabling online decision engines to learn dynamic frugality under fluctuating SLA constraints.
    \item \textbf{Comprehensive Empirical Validation:} Evaluated across 3,712 queries from GSM8K, LogiQA, and CommonsenseQA, AdaptInfer matches top-tier accuracy (81.4\% vs. 82.0\%) while reducing average system latency by 67.2\% (14.8s vs. 45.2s) and token cost by 27.5\% (1,421 vs. 1,959 tokens/query).
    \item \textbf{Full-Stack Real-Time Dashboard:} We deliver an open-source, single-page application dashboard providing live telemetry, strategy racing, confidence bar visualization, and dynamic Chart.js analytics.
\end{enumerate}

\section{System Architecture \& Software Design Patterns}\label{sec2}

The AdaptInfer system architecture is organized into four decoupled layers: the Client Telemetry Layer, the Asynchronous Microservice Layer, the Decision Core, and the Distributed Task Execution Engine. Figure~\ref{fig:arch_diagram} illustrates the end-to-end data flow and structural interaction between components.

\begin{figure*}[htbp]
\centering
\includegraphics[width=0.92\textwidth]{sdp_architecture_diagram.png}
\caption{End-to-End AdaptInfer Distributed Microservices Architecture. Queries submitted via the reactive dashboard are vectorized by Sentence-BERT, evaluated by the dynamic decision core, passed through the Confidence Gate guardrail, and executed across distributed local Ollama worker nodes.}
\label{fig:arch_diagram}
\end{figure*}

\subsection{Asynchronous Microservice Engine (FastAPI)}
To prevent slow LLM generation calls from blocking application threads, the application backend is implemented using **FastAPI** operating on an asynchronous ASGI server (Uvicorn). The backend exposes three primary REST microservice endpoints:
\begin{itemize}
    \item \texttt{POST /api/analyze}: Receives a user question, executes the routing engine, invokes selected strategies asynchronously, calculates correctness and token cost, and returns complete execution telemetry.
    \item \texttt{POST /api/recommend}: Accepts raw prompt text, extracts semantic state vectors, evaluates LinUCB action scores, and returns normalized confidence probabilities for UI bar visualization.
    \item \texttt{GET /api/history}: Serves aggregated historical benchmark metrics across 3,712 queries formatted for real-time Chart.js rendering.
\end{itemize}

\subsection{Distributed Task Execution (Ray Framework \& Local Ollama)}
To isolate model inference from API thread handling, AdaptInfer integrates the **Ray distributed computing framework** \cite{moritz2018ray}. When an action $a_t$ is selected, the task is dispatched to a Ray worker actor process. 

Models (\texttt{Qwen2.5-3B-Instruct} and \texttt{Qwen2.5-7B-Instruct}) are hosted locally using Ollama runtime instances. Local hosting completely eliminates network latency jitter inherent to cloud LLM APIs, ensuring that recorded system latency strictly reflects pure hardware inference processing time.

\subsection{Object-Oriented Design Pattern Implementations}
The codebase adheres strictly to object-oriented software engineering design patterns to maintain high extensibility, maintainability, and testability.

\subsubsection{1. Strategy Pattern}
To ensure all reasoning pathways are interchangeable at runtime, we define an abstract base class \texttt{BaseStrategy} in \texttt{src/strategies/base.py}:
\begin{lstlisting}[language=Python]
class BaseStrategy(ABC):
    @abstractmethod
    def solve(self, question: str) -> StrategyResult:
        """Execute reasoning strategy and return structured results."""
        pass
\end{lstlisting}
Concrete classes (\texttt{DirectStrategy}, \texttt{CoTStrategy}, \texttt{GoTStrategy}, \texttt{ToTStrategy}) inherit from this interface. The router dynamically instantiates and invokes these strategies polymorphically without modifying client caller code.

\subsubsection{2. Chain of Responsibility Pattern (Confidence Gate)}
To mitigate catastrophic errors during early online learning, we implement a \texttt{ConfidenceGate} guardrail. The gate intercepts the action $a_t$ selected by the engine and samples the target model at two distinct temperatures: $T=0.2$ (low temperature, deterministic) and $T=0.8$ (high temperature, exploratory).

If the model genuinely understands the problem, its logit distribution is sharply peaked, producing identical answers ($y_1 \equiv y_2$). If the logit distribution is flat (indicating uncertainty or potential hallucination), $y_1$ and $y_2$ diverge ($y_1 \not\equiv y_2$). Upon detecting divergence, the gate intercepts execution and escalates the query along a defined chain of responsibility:
\begin{equation}
a_{t+1} = \mathcal{E}(a_t)
\end{equation}
where $\mathcal{E}$ maps simple actions to higher-capacity fallback strategies (e.g., $\text{3B\_Direct} \rightarrow \text{7B\_CoT} \rightarrow \text{7B\_ToT}$).

\subsubsection{3. Adapter Pattern}
The low-level HTTP REST interface of the Ollama server is encapsulated within an \texttt{LLMClient} adapter class (\texttt{src/llm/client.py}). The adapter standardizes generation options ($\text{seed}=42$, $\text{temperature}=0.2$, $\text{max\_tokens}=512$, $\text{timeout}=120\text{s}$, $\text{max\_retries}=0$), shielding core reasoning strategies from low-level network transport details.

\section{Contextual Decision Engine \& Dynamic Budget Tracker}\label{sec3}

\subsection{State Representation \& Dynamic Budget Scalar}
Each incoming query $x_t$ is transformed into a dense vector using a pre-trained Sentence-BERT encoder (\texttt{all-MiniLM-L6-v2}) \cite{reimers2019sentence}, yielding a 384-dimensional semantic embedding $\phi(x_t) \in \mathbb{R}^{384}$. 

To enable budget-aware routing, the system maintains a Dynamic Budget Tracker. Let $B_{\max}$ be the maximum token allocation for a user session, and let $C_t$ be the total cumulative tokens consumed up to step $t$. The remaining budget ratio $b_t$ is defined as:
\begin{equation}
b_t = \max\left(0, \frac{B_{\max} - C_t}{B_{\max}}\right) \in [0, 1]
\end{equation}
The final augmented state vector $s_t \in \mathbb{R}^{385}$ is formed by concatenating the budget ratio to the semantic embedding:
\begin{equation}
s_t = \left[ \phi(x_t)^\top, \, b_t \right]^\top
\end{equation}
As $b_t \rightarrow 0$, the decision engine mathematically observes that selecting token-heavy strategies (e.g., 7B\_ToT) incurs severe penalty terms, automatically forcing the system toward lightweight actions.

\subsection{LinUCB Decision Core}
We select the LinUCB algorithm \cite{li2010contextual} to govern action selection across our 6-discrete joint action space:
\begin{equation}
\mathcal{A} = \{ \text{3B\_Direct}, \text{3B\_CoT}, \text{7B\_Direct}, \text{7B\_CoT}, \text{7B\_GoT}, \text{7B\_ToT} \}
\end{equation}
Unlike Deep Neural Network algorithms (e.g., DQN or PPO) which require Markov state transitions, LLM query routing is inherently stateless. LinUCB maintains a linear model for each action $a \in \mathcal{A}$ with parameter vector $\theta_a \in \mathbb{R}^{385}$ and covariance matrix $A_a \in \mathbb{R}^{385 \times 385}$. At step $t$, the agent computes the Upper Confidence Bound for each action:
\begin{equation}
a_t = \arg\max_{a \in \mathcal{A}} \left( \theta_a^\top s_t + \alpha \sqrt{s_t^\top A_a^{-1} s_t} \right)
\end{equation}
where $\alpha > 0$ controls the exploration-exploitation trade-off. Upon observing execution feedback, matrix $A_{a_t}$ and bias vector $b_{a_t}$ are updated in closed form:
\begin{align}
A_{a_t} &\leftarrow A_{a_t} + s_t s_t^\top \\
b_{a_t} &\leftarrow b_{a_t} + r_t s_t \\
\theta_{a_t} &\leftarrow A_{a_t}^{-1} b_{a_t}
\end{align}

\subsection{Multi-Objective Composite Reward Function}
The composite reward $r_t$ evaluates answer correctness while penalizing token cost and response latency:
\begin{equation}
r_t = \lambda_{\text{acc}} \cdot \mathbb{I}(\text{correct}_t) - \lambda_{\text{cost}} \cdot \left(\frac{\text{Tokens}_t}{\text{Tokens}_{\max}}\right) - \lambda_{\text{lat}} \cdot \left(\frac{\text{Latency}_t}{\text{Latency}_{\max}}\right)
\end{equation}
In our implementation, hyperparameters are set to $\lambda_{\text{acc}} = 1.0$, $\lambda_{\text{cost}} = 0.1$, and $\lambda_{\text{lat}} = 0.05$, establishing a strong preference for accuracy while driving token and time frugality.

\section{Empirical Evaluation \& System Results}\label{sec4}

\subsection{Experimental Setup}
The evaluation suite was executed across a clean benchmark stream of $N=3,712$ queries spanning three distinct cognitive domains:
\begin{enumerate}
    \item \textbf{GSM8K (1,319 queries):} Multi-step mathematical word problems requiring quantitative deduction \cite{cobbe2021training}.
    \item \textbf{LogiQA (1,172 queries):} Logical reasoning and reading comprehension passages \cite{liu2020logiqa}.
    \item \textbf{CommonsenseQA (1,221 queries):} General knowledge and physical commonsense reasoning \cite{talmor2019commonsenseqa}.
\end{enumerate}

\subsection{Cost-Accuracy Trade-off Analysis (Pareto Frontier)}
Table~\ref{tab:system_results} details the empirical performance of static routing baselines versus dynamic AdaptInfer routing.

\begin{table}[htbp]
\centering
\caption{Empirical System Performance Breakdown across 3,712 Benchmark Queries. AdaptInfer achieves near statistical parity with static 7B-ToT while slashing latency by 67.2\% and saving 27.5\% in token expenditure.}
\label{tab:system_results}
\begin{tabular}{lcccc}
\toprule
\textbf{Routing Policy} & \textbf{Accuracy (\%)} & \textbf{Avg Tokens / Query} & \textbf{Median Latency (s)} & \textbf{Token Savings} \\
\midrule
3B-Direct & 53.9\% & 108 & 4.4s & -- \\
3B-CoT & 62.5\% & 369 & 26.5s & -- \\
7B-Direct & 57.6\% & 108 & 11.2s & -- \\
7B-CoT & 80.0\% & 379 & 17.4s & -- \\
7B-GoT & 83.5\% & 1,267 & 32.1s & -- \\
7B-ToT (Upper Bound) & \textbf{82.0\%} & \textbf{1,959} & \textbf{45.2s} & Baseline \\
\midrule
\textbf{AdaptInfer (Ours)} & \textbf{81.4\%} & \textbf{1,421} & \textbf{14.8s} & \textbf{-27.5\%} \\
\bottomrule
\end{tabular}
\end{table}

\begin{figure}[htbp]
\centering
\includegraphics[width=0.82\textwidth]{sdp_latency_improvement.png}
\caption{Architectural inference latency optimization: AdaptInfer cuts response latency by 67.2\% relative to static Tree-of-Thoughts.}
\label{fig:latency_graph}
\end{figure}

\begin{figure}[htbp]
\centering
\includegraphics[width=0.82\textwidth]{sdp_token_savings.png}
\caption{Token expenditure reduction: AdaptInfer saves 538 tokens per query (27.5\% net savings) compared to static Tree-of-Thoughts.}
\label{fig:token_graph}
\end{figure}

\begin{figure}[htbp]
\centering
\includegraphics[width=0.82\textwidth]{sdp_system_tradeoff.png}
\caption{Systems Pareto efficiency: Accuracy vs Response Latency across static baselines and AdaptInfer.}
\label{fig:pareto_graph}
\end{figure}

\subsection{Routing Action Distribution}
Table~\ref{tab:distribution} illustrates how the decision engine distributed query traffic across the action space.

\begin{table}[htbp]
\centering
\caption{Final Action Selection Breakdown by AdaptInfer over 3,712 Queries.}
\label{tab:distribution}
\begin{tabular}{lcc}
\toprule
\textbf{Selected Action} & \textbf{Query Count} & \textbf{Selection Percentage} \\
\midrule
7B-Direct & 1,037 & 27.9\% \\
3B-Direct & 742 & 19.9\% \\
7B-CoT & 651 & 17.5\% \\
3B-CoT & 622 & 16.7\% \\
7B-ToT & 358 & 9.6\% \\
7B-GoT & 302 & 8.1\% \\
\bottomrule
\end{tabular}
\end{table}

The engine routed **47.8\% of all traffic** to fast direct prompting (3B/7B Direct) and **34.2\%** to standard Chain-of-Thought, reserving complex Tree/Graph-of-Thoughts strategies for only the top **17.7\% of highly intricate queries**.

\subsection{Domain-Specific Adaptation}
Table~\ref{tab:domain_breakdown} summarizes accuracy across individual datasets. On mathematical reasoning (GSM8K), complex ToT reasoning yields massive accuracy gains (88.4\% vs. 31.2\%). On general commonsense knowledge (CommonsenseQA), direct prompting achieves 82.1\% while ToT achieves 90.5\%, making ToT computationally wasteful. LinUCB automatically learns these domain signatures, avoiding heavy search algorithms on commonsense queries entirely.

\begin{table}[htbp]
\centering
\caption{Domain-Specific Performance Breakdown Across Datasets.}
\label{tab:domain_breakdown}
\begin{tabular}{lccc}
\toprule
\textbf{Routing Policy} & \textbf{GSM8K (Math)} & \textbf{LogiQA (Logic)} & \textbf{CommonsenseQA} \\
\midrule
3B-Direct & 31.2\% & 48.5\% & 82.1\% \\
7B-ToT & 88.4\% & 67.2\% & 90.5\% \\
\textbf{AdaptInfer} & \textbf{87.1\%} & \textbf{66.8\%} & \textbf{90.2\%} \\
\bottomrule
\end{tabular}
\end{table}

\section{Interactive Dashboard \& Live Telemetry}\label{sec5}

To provide complete system explainability, AdaptInfer includes a full-stack, single-page application dashboard running on port 8000.

\subsubsection{Key Dashboard Features:}
\begin{enumerate}
    \item \textbf{Live Strategy Race Panel:} Submits user queries to all 6 strategies simultaneously, displaying live latency, token consumption, and parsed answers.
    \item \textbf{LinUCB Recommendation \& Confidence Bars:} Visualizes decision engine upper confidence bound probability distributions across all 6 actions.
    \item \textbf{AdaptInfer Efficiency Box:} Calculates real-time token savings and latency reductions for the current query against static ToT.
    \item \textbf{Dynamic Agent Performance Summary Table:} Displays live session metrics (Accuracy, Avg Tokens, Median Latency, Avg Reward, Avg LLM Calls) updating dynamically after every query.
\end{enumerate}

\section{Conclusion \& Future Work}\label{sec6}

In this paper, we introduced \textbf{AdaptInfer}, a distributed software architecture for dynamic joint selection of LLM parameter size and reasoning strategy under real-time session budget constraints. By decoupling LLM generation via Ray distributed actors, incorporating a dynamic budget scalar into a LinUCB contextual state vector, and implementing a Confidence-Gated Chain of Responsibility fallback, AdaptInfer bridges the gap between high accuracy and computational efficiency. Across 3,712 benchmark queries, AdaptInfer achieves 81.4\% accuracy while reducing system response latency by 67.2\% and token costs by 27.5\%.

Future work will focus on extending the Ray worker infrastructure to multi-node GPU clusters and incorporating multimodal Vision-Language Models (VLMs) into the dynamic action space.

\backmatter

\bmhead{Declarations}
\begin{itemize}
    \item \textbf{Funding:} The authors declare no external funding was received for this study.
    \item \textbf{Competing Interests:} The authors declare no competing financial or non-financial interests.
    \item \textbf{Availability of Data and Code:} Source code and benchmark evaluation datasets are open-sourced in the project repository.
\end{itemize}

\bibliography{sn-bibliography}

\end{document}
